\section[Framework delegates and scientific completion]{When a completed software call leaves the research unfinished}
\label{sec:exec28-dl}
A causal workflow can finish its control flow without producing an admissible causal answer.
For a collections analyst, the distinction is practical. A reminder-policy analysis may return an
object after its structured-output step failed, repeat a tool after an interrupted response, or
accept a treatment identifier under a weakened constraint. Each event changes what the analyst can
infer from the result. A framework comparison must therefore follow the delegated implementation
of validation, stopping, retry and recovery, rather than infer these properties from a common API.
The following comparison concerns selected, pinned original source paths. It motivates product
contracts and public qualification experiments; it does not report executed vendor integrations.

\subsection{Validation must survive the complete calling path}
LangChain's inspected tool base distinguishes a dictionary schema from a Pydantic model. In the
dictionary branch of the input parser, the supplied input is returned; the Pydantic branch invokes
model validation and handles an injected tool-call identifier. This local branch difference does
not prove that no other layer validates a dictionary-schema request. It does show why the proposed
adapter cannot treat the presence of a schema as evidence that validation occurred at every entry
point. The same source can convert a handled validation or tool exception into an error-status tool
message. The consumer must preserve that status even when the returned content is readable text.
\citep{exec28_dl_langchain}

DSPy's inspected tool helper iterates over provided arguments, conditionally validates a field
against its schema when a qualifying type is present, then uses typed parsing for applicable
arguments. Its synchronous and asynchronous call paths also differ in coroutine handling. Google
GenAI's function-declaration helper builds a schema from a callable signature and type information,
including required fields and definitions; an inspected Vertex conversion branch supplies an
object type for an otherwise untyped union. These are useful mechanisms for translating interfaces.
They do not establish that a variable precedes treatment, that a graph encodes the business question,
or that omitted information is scientifically harmless. Conditional adoption should retain the
original helper while placing a common causal-request validator at the library boundary.
\citep{exec28_dl_dspy,exec28_dl_google}

Consider a request whose treatment identifier must contain a digit. Within the stipulated domain
\(D=\{\texttt{x},\texttt{7}\}\), the intended accepted set is
\(S=\{\texttt{7}\}\). If a conversion relaxes the restriction to any string, the accepted set becomes
\(\widetilde S=D\), so \(S\subsetneq\widetilde S\). The extra value is a false admission under the
original contract. This finite authored witness establishes the logical consequence of weakening;
it is not an observed generation or a measured package failure.~\citep{exec28_dl_reference}

The distinction matters for llama.cpp's inspected schema-to-grammar converter. An unsupported
regular-expression pattern can generate a warning and fall back to any string, whereas an invalid
pattern is recorded as an error. The converter subsequently raises recorded errors and prints
warnings. A separate build/configuration path can use LLGuidance; the warning behavior must not be
generalized to every backend. The selected chat-template path chooses schema-derived grammar or a
supplied grammar, while a separate argument parser can retain text when JSON parsing fails.
Consequently, grammar generation, parse success and validation against the original schema are
three distinct observations. The proposed library should retain conversion diagnostics and the
selected grammar backend, and validate the completed payload against the original request before
admitting an effect. Wholesale porting would inherit these semantic choices and add maintenance
cost without removing the need for this validation.~\citep{exec28_dl_llamacpp}

\subsection{A stop event is a control state, not a scientific certificate}
In LlamaIndex's inspected agent workflow, structured-response generation can raise an exception
that is converted to a warning before the workflow returns a stop event. Its early-stop generation
option can also make another model call to construct a final response. Thus, a terminal event alone
cannot certify a valid structured answer. In LangChain's classic agent, forced stopping returns a
limit response, while the generation option can request a final model answer. The synchronous loop
updates elapsed time between steps; that observation is narrower than a claim that all outstanding
remote operations are cancelled at a wall-clock limit.~\citep{exec28_dl_llama_index,exec28_dl_langchain}

Let \(T\) indicate terminal control flow, \(V\) validation of the scientific envelope, and \(I\)
admission of the requested effect under declared assumptions. The proposed scalar-effect gate is
\(T\land V\land I\). An exact four-case construction admits only the terminal, valid, identified
case. A completed nonidentification result remains a legitimate answer with a different disposition;
a terminal structured-output failure remains a failed representation; an incomplete stream remains
incomplete. The gate specifies responsibility, not the truth of a bank's assumptions. A framework
owns its execution state, the adapter checks representation and preserves errors, and the causal
library supplies the identification argument.~\citep{exec28_dl_reference}

Semantic Kernel's inspected chat agent combines arguments and execution settings before delegating
to a selected provider. Its function wrapper invokes filters, records errors and propagates
exceptions; the streaming path can yield partial results. These hooks are useful integration
locations, but the inspected wrapper does not establish a causal validator. Haystack validates
components, sockets and mandatory inputs; its execution wrapper checks the returned mapping and
output names, and a blocked queue can raise an error. Such checks protect pipeline structure.
They should be retained alongside a separate validation of the target, information clock, graph,
estimand and uncertainty claim. A proposed CLI command and MCP tool should return the same typed
scientific disposition even when their orchestration frameworks differ.
\citep{exec28_dl_semantic_kernel,exec28_dl_haystack}

\subsection{Budgets must count actual attempts and side effects}
LangChain recommends placing retries around the smallest operation that needs them. Its inspected
retry wrapper has configurable exception handling and backoff; the batch path retains successful
items and retries failed original indices. Suppose three independent items \(a,b,c\) run once,
only \(b\) fails, and \(b\) needs two further attempts. Selective retry makes \(1+3+1=5\) item
attempts, whereas repeating the whole three-item batch three times makes nine. At a stipulated
USD 0.25 per item attempt, the costs are USD 1.25 and USD 2.25. These exact accounting examples
are not vendor prices. They also assume that retained successes are valid and that the repeated
operation is safe to retry; neither property follows merely from an exception class.
\citep{exec28_dl_langchain,exec28_dl_reference}

Google GenAI's inspected HTTP path converts configured milliseconds into seconds and supports
request-specific retry settings in addition to client defaults. The inspected server-timeout header
rounds seconds upward. Haystack's visit limit applies to an individual component, while a workflow
can traverse several components. A business-request budget must therefore count HTTP attempts,
tool attempts, generation calls, waiting time and unresolved outcomes explicitly. A framework
iteration cap, a component visit cap and an HTTP timeout are different limits.
\citep{exec28_dl_google,exec28_dl_haystack}

Anthropic's pinned tool-runner documentation describes dispatch by native stop reason. Tool-use
responses trigger tools; paused turns and compaction can require continuation; other stop states
terminate the iterator after yielding the message. Its documented eager streaming mode can execute
a tool before the overall reply finishes. If that reply is interrupted, the tool's effect cannot be
undone merely by discarding the conversation, and the model may request the call again. The same
documentation describes compaction calls outside the ordinary iteration count and nested usage
accounting. Two ordinary turns plus one separate compaction make three calls in an authored counting
example. A top-level counter cannot be assumed to represent the complete cost denominator.
\citep{exec28_dl_anthropic,exec28_dl_reference}

For the reminder-policy workflow, graph inspection and effect calculation can be read-only tools;
dispatching an actual reminder is a different operation. The proposal should assign stable operation
identities, preserve pending results, and require a deduplication or reconciliation policy before
retrying actions with external effects. Cancellation requested and cancellation observed must be
separate states. A failed or unknown outcome remains in the denominator when comparing cost and
reliability. These are proposed architectural requirements, not a claim that any reviewed framework
already supplies end-to-end exactly-once execution.

\subsection{Recovery, native stop reasons and learned programs}
The inspected delegated workflow context describes a process-local store and serialization of
state, event queues, buffers and related execution metadata. Cancellation is delegated for an
external context, while other context states issue warnings. Serialization can support resumption,
but this source path alone cannot establish that a remote action and a local checkpoint commit
atomically. The proposed recovery record must therefore include both local state and unresolved
external-operation identities. Tests should interrupt immediately before and after issuing a tool,
receiving its result and recording the result.~\citep{exec28_dl_workflows}

Native stop reasons also need to survive compatibility adapters. Anthropic's selected type includes
refusal, pause, tool use, output limits and context exhaustion. In vLLM's inspected serving path,
parser selection depends on tool and reasoning settings, strictness, the model and tokenizer.
Its compatibility finish reason can depend on tool mode and generated tool content while retaining
the native stop reason. llama.cpp distinguishes native end-of-sequence, stop-word and limit states
before constructing compatibility fields; the task structure also retains cancellation targets and
parent-child identities. These source representations support a concrete preservation requirement.
They do not independently prove that cancellation has reached every child or that two servers with
similar response fields provide equivalent behavior.~\citep{exec28_dl_anthropic,exec28_dl_vllm,exec28_dl_llamacpp}

Optimization introduces another boundary. DSPy's inspected bootstrap procedure selects traces from
training examples. A variable called validation contains the unbootstrapped training subset, and
success depends on the configured metric; without a metric, the inspected branch accepts success.
The threshold branch also depends on its truth value. These are selected source semantics, not a
claim of a general empirical defect. A causal assistant still needs separately held-out tasks with
independent expected answers, including nonidentification, invalid adjustment and ambiguous graph
cases. Optimizing a demonstration-selection metric cannot by itself establish causal validity on
new banking questions.~\citep{exec28_dl_dspy}

\subsection{A bounded adoption and qualification decision}
The proposal favors thin adapters around selected original clients and orchestration components.
LangChain contributes composable calls and scoped retries; Semantic Kernel contributes provider and
function hooks; Haystack contributes explicit pipeline structure; LlamaIndex and its delegated
workflow package contribute orchestration and recovery representations; DSPy contributes a candidate
optimization layer subject to an independent evaluation contract. The decision is conditional
integration, with scientific validation owned by the causal library. It rejects automatic promotion
of terminal messages into effects, implicit schema weakening, unmeasured retry economics and
unqualified wholesale ports.~\citep{exec28_dl_langchain,exec28_dl_semantic_kernel,exec28_dl_haystack,exec28_dl_llama_index,exec28_dl_workflows,exec28_dl_dspy}

For hosted backends, client source inspection and service qualification remain separate. For local
backends, the exact server build, model, tokenizer, template, parser and grammar configuration are
part of the candidate. Ollama's selected README itself requires a separately running server and a
selected model; its client identity cannot substitute for those artifacts. README migration or
product-focus statements likewise supply leads, not inherited qualification: the Semantic Kernel
README names a successor framework, and the LlamaIndex README names parsing products and benchmarks.
Those separate products require their own primary-source comparison before adoption claims.
\citep{exec28_dl_ollama,exec28_dl_semantic_kernel,exec28_dl_llama_index}

A low-cost public qualification suite can now be specified precisely. It should present an invalid
original-schema value, a handled tool error, a structured-output failure followed by a terminal event,
a partly failed batch, an interrupted side-effecting call, a compaction call outside the ordinary
counter, and a limit stop containing a valid-looking prefix. Each fixture needs an independent expected
scientific disposition, attempt ledger and recovery outcome. Subsequent native trials require exact
artifacts, rights and declared budgets; bank effectiveness remains a later experiment. The comparison
supports these targeted experiments without converting static source inspection into measured merit.
The causal answer is trustworthy only when its meaning survives validation, scheduling, stopping and recovery.
